\chapter{The Class Side}
\label{ch:classside}

\standfirst{Classes are objects, so they have a class of their own, and that
class is where \ct{new} lives. This chapter is the part of Smalltalk that looks
strange for about an hour and then never troubles you again.}

\section{Why there is a class side at all}

If classes are objects---and they are, you can put one in a collection and
send it messages---then a class must be an instance of something, and messages
sent to a class must be looked up somewhere.

They are looked up in the class's \emph{metaclass}. Every class has exactly
one, created with it, and its sole instance is that class.

\begin{st}
3 class                  "SmallInteger        -- the class of 3"
SmallInteger class       "SmallInteger class  -- its metaclass"
SmallInteger class class "Metaclass"
Metaclass class class    "Metaclass -- the regress stops here"
\end{st}

That is the whole of it. \emph{Instance-side} methods are found when you send
to an instance; \emph{class-side} methods are found when you send to the
class.

\section{Using it}

In the Browser, the \ct{instance} and \ct{class} buttons under the class list
switch which side you are editing. Everything else is the same.

The two things class-side methods are actually for:

\subsection{Making instances}

\begin{st}
Account class >> new
    ^super new initialize

Account class >> newWithBalance: anAmount
    ^self new deposit: anAmount; yourself

Point class >> x: anX y: aY
    ^self new setX: anX y: aY
\end{st}

A well-named class-side constructor is better than making callers send
\ct{new} and then a sequence of setters, because it can enforce that the
object is never seen half-built.

\begin{stnote}
\ct{super new} in a class method means "start the lookup of \ct{new} above my
metaclass"---which reaches \ct{Behavior>>new}, the primitive that allocates.
\ct{self new} inside \ct{Account class>>new} would call itself forever.
\end{stnote}

\subsection{Class-side state and initialisation}

A class-side method can read and write class variables, and a class may define
\ct{initialize} on the class side to set them up:

\begin{st}
Account class >> initialize
    InterestRate := 5 / 100
\end{st}

This does not run by itself when you accept it. Send \ct{Account initialize}
once. When a class is loaded from a package, the loader runs class-side
\ct{initialize} for you after everything is in place, which is how the 1983
library's forty-six class initialisers get run at bootstrap.

\begin{stwarn}
Class-side \ct{initialize} runs at build time, and that is deliberate rather
than incidental: initialising eagerly at build is what makes lazy
initialisation---and its check-then-act race across workers---unnecessary. See
Chapter~\ref{ch:contract}.
\end{stwarn}

\section{Class instance variables}

There is a third kind of variable, rarer and occasionally exactly right. A
metaclass can have instance variables of its own; since its sole instance is
the class, each class in a hierarchy gets its own copy.

\begin{center}
\small
\begin{tabular}{@{}lp{0.58\textwidth}@{}}
\toprule
instance variable & one per instance\\
class variable & \emph{one}, shared by the class, its subclasses and all their
  instances\\
class instance variable & one per class, and subclasses get their \emph{own},
  not the same one\\
\bottomrule
\end{tabular}
\end{center}

Use a class instance variable when each subclass should have its own default;
use a class variable when there must be exactly one for the whole hierarchy.

\section{The picture}

\begin{plain}
      Object  ------ instance-of ----->  Object class
        ^                                     ^
     superclass                           superclass
        |                                     |
      Account ------ instance-of ----->  Account class
        ^                                     ^
     superclass                           superclass
        |                                     |
  SavingsAccount -- instance-of --> SavingsAccount class

  and every metaclass is an instance of Metaclass.
\end{plain}

The metaclass hierarchy runs exactly parallel to the class hierarchy. That is
why \ct{super new} in \ct{SavingsAccount class>>new} finds
\ct{Account class>>new}: the metaclasses inherit in the same order the classes
do.

\ct{Object class} inherits from \ct{Class}, which inherits from
\ct{ClassDescription}, then \ct{Behavior}, then \ct{Object}. That is where
\ct{new}, \ct{selectors}, \ct{compile:} and the rest of the reflective
protocol come from, and it is why every class understands them without anybody
writing them down twice.

\section{When you do not need any of this}

Most days. Write instance methods, write a class-side \ct{new} if your class
needs initialising, and let the rest sit there working. The metaclass
machinery is worth understanding because it explains why the system has no
special cases, not because you have to think about it to get work done.
